% Abstract — leads with the order-statistic mechanism

We present the \emph{Estimated Dynamic Equilibrium Model} (EDEM), an
agent-based framework in which supply and demand are realisations of
a coupled stochastic process driven by heterogeneous, error-prone
agent valuations. The framework's central technical contribution is
a generative mechanism for persistent disequilibrium: when
market-clearing prices are repeatedly selected from the upper tail
of noisy bid distributions and fed back into future valuations,
expected prices drift upward even with strictly zero-mean estimation
errors. We derive the bias in closed form for a clean special case
($\sigma(n-1)/(n+1)$ for $n$ iid uniform bids) and show via
simulation that compounding the bias across epochs produces
exponential price growth without any behavioural assumption about
investor optimism. EDEM extends \citet{miller1977}'s
divergence-of-opinion theory to a fully dynamic setting and recovers
Walrasian equilibrium and Miller's static premium as limiting cases.
Eight controlled experiments in the Python ABM framework Mesa, on a
$32 \times 32$ real-estate neighbourhood, demonstrate six
qualitatively distinct regimes --- band-stable, business-cycle,
persistent overshoot, persistent undershoot, runaway bubble, and
constant transition --- all reachable from the same agent ruleset.
Implications for the empirical divergence-of-opinion literature
(which has produced contradictory findings) and for machine-learning
valuation algorithms (which inherit and amplify the order-statistic
bias) are discussed.
